\documentclass[11pt]{article}
\usepackage[margin=1in]{geometry}
\usepackage{amsmath,amssymb,booktabs,array,longtable,enumitem}
\usepackage[hidelinks]{hyperref}

\newcommand{\E}{\mathbb{E}}
\newcommand{\Var}{\operatorname{Var}}
\newcommand{\Cov}{\operatorname{Cov}}
\newcommand{\R}{\mathbb{R}}
\newcommand{\MSE}{\operatorname{MSE}}
\newcommand{\supp}{\operatorname{supp}}
\newcommand{\plim}{\operatorname{plim}}

\title{Exposure-Map Feedback and AI-Validation Design:\\
A July 31 Continuation Note}
\author{}
\date{July 31, 2026}

\begin{document}
\maketitle

\section*{Purpose and source convention}

This is a new, separate continuation note. It does not modify the July 27 or
July 30 exposure drafts, the feedback documents, or the earlier AI-project
note. It answers four questions:
\begin{enumerate}[leftmargin=2em]
    \item Does the injectivity problem for a scalar exposure index require a
    continuous parameter, and how does it depend on finite, countable, or
    uncountable support?
    \item Does replacing \(Y_i(W)\) by \(Y_i\) solve the graph-HAC centering
    problem?
    \item Which feedback items were actually addressed between the July 27 and
    July 30 exposure drafts?
    \item How does AI-project MSE depend on the number of validated cells, what
    is the resulting threshold rule, and when must validation sites be sampled
    rather than chosen for convenience?
\end{enumerate}

The exposure sources are the July 27 feedback, the July 29 proof-review note,
the July 27 draft, the July 30 redline, Gao, Harshaw, S\"avje, and Wang (2026),
and the earlier June exposure reviews in the project record. The AI sources are
the project proposal and
\texttt{ai\_project\_three\_papers\_whiteboard\_note.tex}.
Statements labeled ``derived'' are derived below.

\section{Exposure-map feedback}

\subsection*{Notation}

\begin{longtable}{>{\raggedright\arraybackslash}p{0.23\linewidth}
                  >{\raggedright\arraybackslash}p{0.69\linewidth}}
\toprule
Symbol & Meaning \\
\midrule
\endfirsthead
\toprule
Symbol & Meaning \\
\midrule
\endhead
\(n,N_n\) & Sample index or number of units, and number of analysis units in
sample \(n\). \\
\(i,v,m\) & Unit index, village index, and number of included annuli. \\
\(D_n\) & Known assignment law in sample \(n\). \\
\(\mathcal A_n\) & Ambient assignment space. \\
\(\mathcal W_n=\supp(D_n)\) & Support of assignments that the design can
generate. \\
\(W\) & Random assignment drawn from \(D_n\). \\
\(Y_i(\cdot)\) & Fixed potential-outcome schedule for unit \(i\). \\
\(Y_i=Y_i(W)\) & Realized-outcome random variable under the design. \\
\(\widetilde Y_i\) & Outcome function of the sufficient exposure at the truth. \\
\(W',y_i^{\rm obs}\) & Fresh placebo assignment and the fixed observed outcome
used in the comparison thought experiment. \\
\(A_v(W)\in\R^m\) & Vector of the first \(m\) annular assignment summaries for
village \(v\). \\
\(S_v=\supp\{A_v(W)\}\) & Support of the annular vector under \(D_n\). \\
\(a,a',h\) & Two distinct support points and a generic support difference. \\
\(\theta\in\Theta\) & Candidate exposure-map parameter. \\
\(\theta_m\in\Theta_m\) & Annular-index weight vector, normalized by
\(\mathbf 1'\theta_m=1\). \\
\(g_v^{(m)}(W;\theta_m)\) & Scalar annular index
\(\theta_m'A_v(W)\). \\
\(H_m,\mathcal B(S_v)\) & Weight-normalization hyperplane and collision set. \\
\(\psi(W)\) & Design-side function used in a moment. \\
\(\phi(Y_i)\) & Outcome transformation used in a moment. \\
\(R_{i,\theta}\psi(W)\) & Exact design residual
\(\psi(W)-\E[\psi(W)\mid g_i(W;\theta)]\). \\
\(\Psi_i(\theta)\) & Unit score
\(\phi(Y_i)R_{i,\theta}\psi(W)\), or a vector of such scores. \\
\(Z_{i,n}(\theta)\) & Infeasibly centered score
\(\Psi_i(\theta)-\E[\Psi_i(\theta)]\). \\
\(\theta_0,\widehat\theta_n\) & True exposure parameter and its estimator. \\
\(\widehat\Omega_n^u(\theta)\) & Feasible graph-HAC estimator formed from
uncentered scores. \\
\(\Omega\) & Limiting covariance matrix at \(\theta_0\). \\
\(M,M_n\) & Fixed and growing numbers of moment functions. \\
\(b_n\) & Maximal affinity-set size. \\
\(V_M^\star,V^\star\) & Finite-dictionary and limiting oracle variances. \\
\(B_n,A,\widehat A,K_\beta,V_\beta\) & Stage-2 centering-bias matrix,
population and estimated derivative matrices, selector for \(\beta\), and the
asymptotic variance of \(\widehat\beta\). \\
\bottomrule
\end{longtable}

\subsection*{1.1 The exact injectivity condition}

The scalar annular index is
\[
    g_v^{(m)}(W;\theta_m)=\theta_m'A_v(W),
    \qquad
    \mathbf 1'\theta_m=1.
\]
For a fixed support \(S_v\), define the collision set
\[
    \mathcal B(S_v)
    =
    \bigcup_{\substack{a,a'\in S_v\\a\neq a'}}
    \left\{
       \theta_m\in\Theta_m:
       \theta_m'(a-a')=0
    \right\}.
\]
The index is one-to-one on \(S_v\) exactly when
\[
    \theta_m\notin\mathcal B(S_v).
\]
Equivalently, using the difference set
\[
    S_v-S_v=\{a-a':a,a'\in S_v\},
\]
injectivity is equivalent to
\[
    \ker(\theta_m')\cap
    \{(S_v-S_v)\setminus\{0\}\}
    =
    \varnothing.
\]
This is the general condition. Cardinality is useful for diagnosing it, but
cardinality alone is not the condition.

\subsection*{1.2 Finite support}

Suppose \(S_v\) is finite and \(\Theta_m\) has nonempty relative interior in
the normalization hyperplane
\[
    H_m=\{\theta_m:\mathbf 1'\theta_m=1\}.
\]
For each pair \(a\neq a'\), the collision equation
\[
    \theta_m'(a-a')=0
\]
cuts out a lower-dimensional hyperplane in \(H_m\), or is empty there. Since
there are finitely many pairs, \(\mathcal B(S_v)\) is a finite union of
lower-dimensional sets. Therefore almost every \(\theta_m\), with respect to
Lebesgue measure on \(H_m\), makes the scalar index injective on \(S_v\).
This is the finite-support result in the July 29 proof-review note.

Continuity of the map \(\theta_m\mapsto g_v^{(m)}\) is not what causes the
problem. The relevant feature is that the admissible parameter set contains a
weight vector outside the collision hyperplanes. A finite or discrete
parameter grid can also contain an injective weight vector; it simply no
longer supports an ``almost every weight'' statement. Conversely, a continuous
parameter set can be deliberately restricted to collision-producing weights,
such as equal annular weights.

\subsection*{1.3 Countably infinite support}

The finite-support conclusion extends to countably infinite support. If \(S_v\)
is countable, then the set of distinct pairs is countable, so
\[
    \mathcal B(S_v)
\]
is a countable union of lower-dimensional hyperplanes. It still has Lebesgue
measure zero in a full-dimensional parameter region. Thus, for almost every
continuously valued weight vector, the scalar index is injective on a
countably infinite support.

A simple example is \(S_v\subset\mathbb Z^m\). A vector whose coordinates are
rationally independent can distinguish all integer vectors in \(S_v\).
Therefore moving from finite to countably infinite support does not by itself
restore exact residual variation.

The same point survives a triangular array. If every \(S_{v,n}\) is finite,
then
\[
    \bigcup_{n=1}^{\infty}\mathcal B(S_{v,n})
\]
is a countable union of hyperplanes. A generic fixed weight vector can be
injective for every \(n\) simultaneously. Merely allowing the assignment
support to grow with the sample is not a remedy.

\subsection*{1.4 Uncountable support}

Uncountability alone is not sufficient to prevent injectivity. For example,
the uncountable set
\[
    S=\{(t,0,\ldots,0):t\in[0,1]\}
\]
is mapped injectively by every \(\theta\) with \(\theta_1\neq0\).

Geometry supplies the useful distinction. If \(m>1\) and \(S_v\) contains an
open subset of \(\R^m\), no scalar linear map can be injective on \(S_v\).
Every nonzero \(\theta_m\) has an \((m-1)\)-dimensional kernel. For an interior
point \(a\in S_v\), a sufficiently small nonzero \(h\in\ker(\theta_m')\)
satisfies both \(a\in S_v\) and \(a+h\in S_v\), while
\[
    \theta_m'a=\theta_m'(a+h).
\]
The same conclusion holds under other full-dimensional support conditions.
But an uncountable support lying on a one-dimensional curve can remain
injectively encodable. The correct question is whether the support has
nontrivial differences in \(\ker(\theta_m')\), not simply whether it is
uncountable.

\subsection*{1.5 Why this does not doom every exposure map}

Suppose \(\psi(W)=h\{A_v(W)\}\). If
\(\theta_m'A_v(W)\) is injective on \(S_v\), then \(A_v(W)\), and hence
\(\psi(W)\), is measurable with respect to the scalar index. Therefore
\[
    R_{v,\theta_m}\psi(W)=0
\]
almost surely. This makes the corresponding exact moment uninformative.

This conclusion is important but narrower than ``every exposure map is
doomed.'' Residual variation remains when, for example:
\begin{enumerate}[leftmargin=2em]
    \item the exposure deliberately coarsens assignment, such as a treated
    count with many assignments producing the same count;
    \item symmetry or equal-weight restrictions force collisions;
    \item the design function \(\psi(W)\) depends on assignment features not
    recoverable from the exposure;
    \item the support is full-dimensional relative to the scalar index; or
    \item the researcher bins or otherwise coarsens the exposure as part of the
    maintained estimand, rather than using approximation error to imitate
    coarsening.
\end{enumerate}

This degeneracy is also distinct from the Gao--Harshaw--S\"avje--Wang
impossibility theorem. Their theorem says no test is uniformly informative
against every bounded potential-outcome schedule in an unrestricted larger
exposure model. It does not say that every test has zero power against every
specific structured alternative. Park's moments may be useful against
restricted alternatives visible through the selected moments, provided the
exact residuals are nondegenerate. Injectivity is an additional failure mode:
for an injective candidate, Park's selected exact moment itself can collapse
before the minimax issue is reached. (Gao et al., Theorem 3.1 and Section 4,
PDF pp.~10--12.)

\subsection*{1.6 Application implication}

For every estimated annular support and weight vector, the application should
report:
\begin{enumerate}[leftmargin=2em]
    \item the probability that distinct annular vectors have the same index;
    \item the design variance of each exact residual;
    \item projection-approximation error relative to exact residual variance;
    and
    \item the smallest singular value of the Stage-1 Jacobian.
\end{enumerate}
If exact collisions are numerically delicate, the paper should state the
tolerance used and show a tolerance path. A quadratic smoother cannot repair
an exactly zero residual. If its approximation error is first-order
negligible, the estimated moment inherits the degeneracy. If the error is
first-order relevant, the moment is driven by approximation error.

\subsection*{1.7 Replacing \(Y_i(W)\) by \(Y_i\) does not fix feedback 0.3}

Under the design-based convention, \(Y_i(\cdot)\) is a fixed schedule and
\[
    Y_i:=Y_i(W)
\]
is the realized-outcome random variable. These are different notational
objects, but \(Y_i\) is still a function of random assignment. Thus
\[
    \E_D[\Psi_i(\theta)]
    =
    \E_D\!\left[
       \phi\{Y_i(W)\}R_{i,\theta}\psi(W)
    \right]
\]
can be written explicitly as
\[
    \sum_{w\in\mathcal W_n}
       D_n(w)\,
       \phi\{Y_i(w)\}
       R_{i,\theta}\psi(w)
\]
for a finite design. Away from \(\theta_0\), this expression generally depends
on potential outcomes \(Y_i(w)\) at assignments that were not realized.
Writing \(\phi(Y_i)\) instead of \(\phi\{Y_i(W)\}\) shortens the notation; it
does not change this expectation or make it computable.

One could instead hold the observed value \(y_i^{\rm obs}\) fixed while drawing
a fresh placebo assignment \(W'\). Then
\[
    \E_{W'}\!\left[
       \phi(y_i^{\rm obs})R_{i,\theta}\psi(W')
    \right]
    =
    \phi(y_i^{\rm obs})
    \E_{W'}[R_{i,\theta}\psi(W')]
    =
    0
\]
for every \(\theta\). But this is a different probability experiment. It is
not the design expectation of the original score, because it artificially
holds the outcome fixed as assignment changes. It therefore cannot be used to
justify the claimed centering of the design-based score.

At the truth, exposure sufficiency does solve the centering problem:
\[
    Y_i=\widetilde Y_i\{g_i(W;\theta_0)\}.
\]
Then
\begin{align*}
    \E_D[\Psi_i(\theta_0)]
    &=
    \E_D\!\left[
       \phi\!\left(\widetilde Y_i\{g_i(W;\theta_0)\}\right)
       R_{i,\theta_0}\psi(W)
    \right]\\
    &=
    \E_D\!\left[
       \phi\!\left(\widetilde Y_i\{g_i(W;\theta_0)\}\right)
       \E_D\{R_{i,\theta_0}\psi(W)\mid g_i(W;\theta_0)\}
    \right]\\
    &=0.
\end{align*}
Therefore the feasible graph-HAC estimator may use uncentered scores at
\(\theta_0\). Evaluating at \(\widehat\theta_n\) still requires the separate
plug-in condition
\[
    \left\|
       \widehat\Omega_n^u(\widehat\theta_n)
       -
       \widehat\Omega_n^u(\theta_0)
    \right\|_{\rm op}
    =
    o_p(1),
\]
together with consistency at the truth,
\[
    \left\|
       \widehat\Omega_n^u(\theta_0)-\Omega
    \right\|_{\rm op}
    =
    o_p(1).
\]
For a smooth parameter this can follow from a local Lipschitz or derivative
bound and the rate of \(\widehat\theta_n\). For a ring it needs a local
shell-oscillation condition. Consistency of \(\widehat\theta_n\) alone is not
enough.

\subsection*{1.8 July 27 versus July 30: July 27 feedback ledger}

This comparison is restricted to items 0.1--0.5 in the July 27 feedback
document.

\begin{longtable}{>{\raggedright\arraybackslash}p{0.09\linewidth}
                  >{\raggedright\arraybackslash}p{0.39\linewidth}
                  >{\raggedright\arraybackslash}p{0.13\linewidth}
                  >{\raggedright\arraybackslash}p{0.29\linewidth}}
\toprule
Item & Requested correction & Status in July 30 & Evidence \\
\midrule
\endfirsthead
\toprule
Item & Requested correction & Status in July 30 & Evidence \\
\midrule
\endhead
0.1 & Diagnose injectivity of the annular index; report collisions, exact
residual variance, approximation error, and Jacobian strength. &
Not addressed &
The annular index remains on p.~30, but the application reports none of the
four requested diagnostics. \\
\addlinespace
0.2 & State that smoothing estimates the exact projection and cannot create
design variation when the exact residual is zero. &
Not addressed &
Remark 2.8 retains the sentence saying pooling over neighboring exposure
values is the response to degeneracy (p.~13). \\
\addlinespace
0.3 & Use feasible uncentered scores at the truth and add a local plug-in
continuity condition for graph-HAC. &
Not addressed &
Assumption 3.9 still says the design centering is computable and invokes only
consistency of \(\widehat\theta_n\) (p.~21). \\
\addlinespace
0.4 & Do not present \(M_n^2b_n/N_n\to0\) as alone proving growing-sieve
attainment. &
Addressed &
The July 30 draft replaces the assumed growing-dimensional expansion with
pointwise fixed-\(M\) feasible theory and a diagonal argument establishing the
existence of a sufficiently slow deterministic \(M_n\). It explicitly says the
rate does not validate every prespecified sequence (pp.~22 and 71--73). \\
\addlinespace
0.5 & Replace expectation ordering by a probability-limit result and require
\(\widehat A\to_p A\). &
Not addressed &
Remark D.5 still begins from \(\E[\widehat\Omega_\Phi]\) and concludes
asymptotic conservativeness without proving the needed probability limit
(pp.~79--80). \\
\bottomrule
\end{longtable}

The July 30 file is a redline rather than a clean replacement. Its main
substantive theoretical improvement is the new slow-sieve theorem. It also
reorganizes and expands the applications. Those edits do not supply the
annular residual diagnostics or the other requested corrections.

\section{AI-project MSE and validation design}

\subsection*{Notation}

\begin{longtable}{>{\raggedright\arraybackslash}p{0.23\linewidth}
                  >{\raggedright\arraybackslash}p{0.69\linewidth}}
\toprule
Symbol & Meaning \\
\midrule
\endfirsthead
\toprule
Symbol & Meaning \\
\midrule
\endhead
\(j=1,\ldots,J\) & Occupation-task cell and total number of target cells. \\
\(\theta_j\) & True productivity gain in cell \(j\). \\
\(p_j,\bar\theta\) & Cell weight and aggregate
\(\bar\theta=\sum_jp_j\theta_j\). \\
\(S_j\) & Broad score observed in every cell. \\
\(\mu,\tau^2,a,b,\sigma_S^2\) & Gaussian score-channel parameters. \\
\(\psi\) & Parameter vector
\((\mu,\tau^2,a,b,\sigma_S^2)\). \\
\(\eta_j\) & Score error. \\
\(R_j,\pi_j,\pi_{jk}\) & Validation indicator, first-order inclusion
probability, and joint inclusion probability. \\
\(m\) & Number of validated cells. \\
\(\widehat v_j,\varepsilon_j,s_j^2\) & Experimental validation estimate,
experimental error, and its design variance. \\
\(h(S_j;\psi)\) & Oracle posterior mean based on the score; this is \(m_j\) in
the earlier AI note. \\
\(\widehat\theta_j\) & Feasible held-out-cell prediction
\(h(S_j;\widehat\psi)\). \\
\(D_S,x_j\) & Score variance \(b^2\tau^2+\sigma_S^2\) and centered score
\(S_j-a-b\mu\). \\
\(\omega^2\) & Oracle posterior variance given only the score. \\
\(\omega_v^2\) & Oracle posterior variance given the score and a validation
estimate with common variance \(s^2\). \\
\(\lambda\) & Score reliability ratio
\(b^2\tau^2/(b^2\tau^2+\sigma_S^2)\). \\
\(\widehat\psi,V_\psi\) & Estimated channel parameter and the covariance in
\(\sqrt m(\widehat\psi-\psi_0)\Rightarrow N(0,V_\psi)\). \\
\(\dot h(S_j;\psi_0)\) & Gradient of the posterior mean with respect to
\(\psi\). \\
\(\Lambda,\Lambda_{\bar\theta}\) & Delta-method sensitivity constants for a
held-out cell and the aggregate. \\
\(I_\psi\{s^2(m)\}\) & Per-cell Fisher information when experimental
precision depends on \(m\). \\
\(T_{\rm MSE},\epsilon,r,\kappa\) & Total-MSE target, excess-MSE tolerance,
RMSE target, and marginal-value tolerance. \\
\(\widehat{\bar\theta}_{\rm EB},
\widehat{\bar\theta}_{\rm PPI}\) & Empirical-Bayes aggregate and design-based
difference estimator. \\
\(g_j=g(S_j)\) & Fixed calibrated prediction used in the design-based
difference estimator. \\
\(e_j,S_e^2\) & Cell residual \(\theta_j-g_j\) and its finite-population
variance. \\
\(S_e,\sigma_v\) & Square roots of \(S_e^2\) and \(\sigma_v^2\). \\
\(\bar s^2\) & Population average \(J^{-1}\sum_js_j^2\). \\
\(B,n_j,\sigma_v^2\) & Total worker-task experimental budget, allocation to
cell \(j\), and per-observation validation variance. \\
\(c\) & Per-site budget overhead, measured in worker-task observation
equivalents. \\
\(d,D\) & Enterprise/site index and number of independently sampled sites. \\
\(Q_d,q_d\) & Site inclusion indicator and probability. \\
\(r_{j\mid d}\) & Conditional cell-inclusion probability within site \(d\). \\
\(X_d,X_j\) & Pre-treatment site and cell covariates used for stratification
or transport. \\
\(\rho\) & Intraclass correlation among validation information from cells in
the same site. \\
\(m_{\rm eff}\) & Approximate effective number of independent validation
cells. \\
\(\bar c=m/D\) & Average number of validated cells per site. \\
\bottomrule
\end{longtable}

\subsection*{2.1 The Gaussian channel and its oracle floor}

The project model is
\[
    \theta_j\sim N(\mu,\tau^2),
\]
\[
    S_j=a+b\theta_j+\eta_j,
    \qquad
    \eta_j\sim N(0,\sigma_S^2).
\]
Write
\[
    D_S=b^2\tau^2+\sigma_S^2,
    \qquad
    x_j=S_j-a-b\mu.
\]
The oracle posterior mean for an unvalidated cell is
\[
    h(S_j;\psi)
    =
    \mu+\frac{b\tau^2}{D_S}x_j,
\]
and the irreducible posterior variance is
\[
    \omega^2
    =
    \frac{\tau^2\sigma_S^2}{D_S}.
\]
Even if the channel parameters were known exactly, the held-out-cell MSE would
not fall below \(\omega^2\). The \(m^{-1}\) statement in the proposal must
therefore refer to \emph{excess risk over the oracle floor}, not total MSE.

\subsection*{2.2 Held-out-cell MSE as a function of \(m\)}

Suppose \(\widehat\psi\) is learned from \(m\) randomly validated cells and is
independent of the held-out cell being evaluated. Assume
\[
    \sqrt m(\widehat\psi-\psi_0)
    \Rightarrow
    N(0,V_\psi).
\]
The prediction is
\[
    \widehat\theta_j=h(S_j;\widehat\psi).
\]
Add and subtract the oracle posterior mean:
\[
    \widehat\theta_j-\theta_j
    =
    \{h(S_j;\psi_0)-\theta_j\}
    +
    \{h(S_j;\widehat\psi)-h(S_j;\psi_0)\}.
\]
The cross term is zero because
\[
    \E[\theta_j-h(S_j;\psi_0)\mid S_j]=0
\]
and the validation sample is independent of the held-out cell. Therefore
\begin{align*}
    \MSE_{\rm held}(m)
    &=
    \omega^2
    +
    \E\!\left[
      \{h(S_j;\widehat\psi)-h(S_j;\psi_0)\}^2
    \right]\\
    &=
    \omega^2+\frac{\Lambda}{m}+o(m^{-1}),
\end{align*}
where
\[
    \Lambda
    =
    \E\!\left[
      \dot h(S_j;\psi_0)'V_\psi
      \dot h(S_j;\psi_0)
    \right].
\]
This is the proposal's claimed expansion with the sensitivity constant made
explicit as a delta-method quadratic form. It is a derived expansion under the
displayed asymptotic-normality and sample-separation conditions; it is not
proved by the five-page proposal.

For completeness, order the parameters as
\[
    \psi=(\mu,\tau^2,a,b,\sigma_S^2).
\]
Then
\[
\dot h(S_j;\psi_0)
=
\begin{pmatrix}
\displaystyle \frac{\sigma_S^2}{D_S}\\[0.8em]
\displaystyle \frac{b\sigma_S^2x_j}{D_S^2}\\[0.8em]
\displaystyle -\frac{b\tau^2}{D_S}\\[0.8em]
\displaystyle
\frac{\tau^2(\sigma_S^2-b^2\tau^2)x_j}{D_S^2}
-\frac{b\tau^2\mu}{D_S}\\[0.8em]
\displaystyle -\frac{b\tau^2x_j}{D_S^2}
\end{pmatrix}.
\]
Thus \(\Lambda\) is computable once the validation design determines
\(V_\psi\).

A transparent special case illustrates the rate. Suppose all channel
parameters except \(\mu\) are known and
\[
    \widehat\mu=\frac1m\sum_{j:R_j=1}\widehat v_j,
    \qquad
    s_j^2=s^2.
\]
Then
\[
    \Var(\widehat\mu)=\frac{\tau^2+s^2}{m},
\]
and
\[
    \frac{\partial h}{\partial\mu}
    =
    1-\lambda
    =
    \frac{\sigma_S^2}{D_S}.
\]
Hence
\[
    \MSE_{\rm held}(m)
    =
    \omega^2
    +
    \frac{(1-\lambda)^2(\tau^2+s^2)}{m}.
\]

\subsection*{2.3 MSE thresholds and the meaning of a ``knee''}

Let \(T_{\rm MSE}\) be a desired total held-out-cell MSE. Since the oracle floor
is \(\omega^2\), the target is feasible only if
\[
    T_{\rm MSE}>\omega^2.
\]
Ignoring the lower-order term, the required number of validated cells is
\[
    m_{\rm MSE}
    \approx
    \left\lceil
      \frac{\Lambda}{T_{\rm MSE}-\omega^2}
    \right\rceil.
\]
If the goal is to come within \(\epsilon\) MSE units of the oracle floor, then
\[
    m_\epsilon
    \approx
    \left\lceil\frac{\Lambda}{\epsilon}\right\rceil.
\]
For an RMSE target \(r>\omega\),
\[
    m_{\rm RMSE}
    \approx
    \left\lceil
      \frac{\Lambda}{r^2-\omega^2}
    \right\rceil.
\]
The marginal MSE reduction from one additional cell is approximately
\[
    \MSE(m)-\MSE(m+1)
    =
    \frac{\Lambda}{m(m+1)}
    \approx
    \frac{\Lambda}{m^2}.
\]
A ``knee'' can therefore be defined by a marginal-value tolerance
\(\kappa>0\):
\[
    m_{\rm knee}\approx\sqrt{\frac{\Lambda}{\kappa}}.
\]
The proposal's numerical value \(m\approx40\) cannot be justified from the
rate alone. It requires an empirical estimate or simulation calibration of
\(\omega^2\), \(V_\psi\), and hence \(\Lambda\), together with a stated
accuracy or marginal-value threshold.

\subsection*{2.4 Full-map risk when validated cells are also evaluated}

If a validated cell has
\[
    \widehat v_j=\theta_j+\varepsilon_j,
    \qquad
    \Var(\varepsilon_j)=s^2,
\]
its oracle posterior variance is
\[
    \omega_v^2
    =
    \left(
      \frac1{\tau^2}
      \frac{b^2}{\sigma_S^2}
      \frac1{s^2}
    \right)^{-1}.
\]
If \(m\) of \(J\) cells are randomly validated, the oracle average map risk is
\[
    R_{\rm oracle,map}(m)
    =
    \left(1-\frac mJ\right)\omega^2
    +
    \frac mJ\omega_v^2.
\]
With cross-fitting to separate channel estimation from cell prediction, a
first-order feasible-risk approximation is
\[
    R_{\rm map}(m)
    =
    \left(1-\frac mJ\right)\omega^2
    +
    \frac mJ\omega_v^2
    +
    \frac{\overline\Lambda(m)}{m}
    +
    o(m^{-1}).
\]
The held-out-cell curve in the proposal instead excludes the directly
validated cells, so its oracle floor remains \(\omega^2\). These two risk
curves should not be conflated.

\subsection*{2.5 Aggregate empirical-Bayes MSE}

Define
\[
    \widehat{\bar\theta}_{\rm EB}
    =
    \sum_{j=1}^Jp_jh(S_j;\widehat\psi).
\]
Conditional on the scores and under independent Gaussian cell residuals,
\[
    \Var\!\left[
       \sum_jp_j\{h(S_j;\psi_0)-\theta_j\}
       \,\middle|\, S_1,\ldots,S_J
    \right]
    =
    \omega^2\sum_jp_j^2.
\]
Let
\[
    \bar d_p=\sum_{j=1}^Jp_j\dot h(S_j;\psi_0).
\]
Then
\[
    \MSE(\widehat{\bar\theta}_{\rm EB})
    =
    \omega^2\sum_jp_j^2
    +
    \frac{\Lambda_{\bar\theta}}{m}
    +
    o(m^{-1}),
\]
where
\[
    \Lambda_{\bar\theta}
    =
    \bar d_p'V_\psi\bar d_p.
\]
With equal weights, the oracle term is \(\omega^2/J\). If \(J\) is large
relative to \(m\), the channel-estimation term dominates, so aggregate RMSE is
of order \(m^{-1/2}\), as asserted in the proposal.

\subsection*{2.6 Exact design-based MSE for the aggregate}

The aggregate does not require the Gaussian channel. Let
\[
    g_j=g(S_j),
    \qquad
    e_j=\theta_j-g_j,
\]
with \(g_j\) fixed relative to validation selection and experimental noise.
Under a simple random sample without replacement of \(m\) cells,
\[
    \widehat{\bar\theta}_{\rm PPI}
    =
    \frac1J\sum_{j=1}^Jg_j
    +
    \frac1m\sum_{j:R_j=1}
       (\widehat v_j-g_j).
\]
Let
\[
    S_e^2
    =
    \frac1{J-1}
    \sum_{j=1}^J(e_j-\bar e)^2,
    \qquad
    \bar s^2=\frac1J\sum_{j=1}^Js_j^2.
\]
Conditional on the finite cell population, with independent experimental
errors,
\[
    \MSE_{\rm PPI}(m)
    =
    \frac1m\left(1-\frac mJ\right)S_e^2
    +
    \frac{\bar s^2}{m}.
\]
Equivalently,
\[
    \MSE_{\rm PPI}(m)
    =
    \frac{S_e^2+\bar s^2}{m}
    -
    \frac{S_e^2}{J}.
\]
The first term is the error from sampling only \(m\) of the finite population
of residuals. The second is experimental estimation noise.

For unequal weights and a general known sampling design, the estimator is
\[
    \widehat{\bar\theta}_{\rm PPI}
    =
    \sum_jp_jg_j
    +
    \sum_j\frac{R_jp_j}{\pi_j}
       (\widehat v_j-g_j).
\]
Ignoring cross-cell experimental-error covariance, its design variance is
\[
    \sum_{j=1}^J\sum_{k=1}^J
    \frac{\pi_{jk}-\pi_j\pi_k}{\pi_j\pi_k}
    p_jp_ke_je_k
    +
    \sum_{j=1}^J\frac{p_j^2}{\pi_j}s_j^2.
\]
This formula makes the role of site clustering and unequal validation
probabilities explicit.

\subsection*{2.7 Fixed precision per cell versus a fixed total budget}

\paragraph{Fixed precision per validated cell.}
If \(s_j^2=s^2\) does not change with \(m\), then
\[
    \MSE_{\rm PPI}(m)
    =
    \frac1m\left(1-\frac mJ\right)S_e^2+\frac{s^2}{m}.
\]
Both the sampling component and the experimental-noise component fall at rate
\(m^{-1}\).

\paragraph{Fixed total worker-task budget.}
Suppose the total number of experimental worker-task observations is \(B\),
allocated evenly across \(m\) cells:
\[
    n_j=\frac Bm.
\]
If the per-observation variance is \(\sigma_v^2\), then
\[
    s_j^2
    =
    \frac{\sigma_v^2}{n_j}
    =
    \frac{\sigma_v^2m}{B}.
\]
The design-based aggregate MSE becomes
\[
    \MSE_{\rm PPI}(m;B)
    =
    S_e^2\left(\frac1m-\frac1J\right)
    +
    \frac{\sigma_v^2}{B}.
\]
Increasing the number of sites or cells reduces the population-coverage
component, but it does not reduce the total within-cell experimental-noise
component when \(B\) is fixed. Without site overhead, using more representative
cells weakly improves this aggregate estimator.

If each new site consumes an overhead of \(c\) observation-equivalents, the
usable within-cell budget is \(B-cm\), and
\[
    \MSE_{\rm PPI}(m;B,c)
    =
    S_e^2\left(\frac1m-\frac1J\right)
    +
    \frac{\sigma_v^2}{B-cm},
    \qquad
    m<\frac Bc.
\]
The continuous interior optimum, before imposing integer and \(m\leq J\)
constraints, is
\[
    m^\star
    =
    \frac{S_eB}{cS_e+\sigma_v\sqrt c},
\]
where \(S_e=\sqrt{S_e^2}\) and
\(\sigma_v=\sqrt{\sigma_v^2}\).

For empirical Bayes, fixed total budget also changes \(V_\psi\). The general
first-order formula becomes
\[
    \MSE_{\rm held}(m;B)
    =
    \omega^2
    +
    \frac1m
    \E\!\left[
       \dot h'
       I_\psi\{s^2(m)\}^{-1}
       \dot h
    \right]
    +
    o(m^{-1}),
\]
where \(I_\psi\{s^2(m)\}\) is per-cell Fisher information and
\[
    s^2(m)=\frac{\sigma_v^2m}{B}.
\]
Because per-cell measurements become noisier as \(m\) grows, the constant is
not fixed. For example, information about \(\mu\) from validation outcomes is
proportional to
\[
    \frac{m}{\tau^2+\sigma_v^2m/B}
    \longrightarrow
    \frac{B}{\sigma_v^2}.
\]
Thus pure \(m^{-1}\) decay cannot be extrapolated indefinitely under a fixed
worker-level budget. The number of cells and the precision per cell must be
chosen jointly.

\subsection*{2.8 Must the project choose which sites receive RCTs?}

Yes, if the target is an economy-wide cell map rather than the causal effect in
the firms that happen to participate. There are two separate randomizations:
\begin{enumerate}[leftmargin=2em]
    \item \emph{Site or cell sampling} determines whether the validated cells
    represent the target population.
    \item \emph{AI-access randomization within a chosen site} identifies the
    causal effect for that chosen site and cell.
\end{enumerate}
The second does not supply the first. Randomizing AI access, model tier, or
``technology per se'' within a convenience partner produces internal validity
for that partner. It does not make the partner's occupation-task effects
representative of the national target.

Unconditional
\[
    R_j\perp\theta_j
\]
is sufficient but not necessary. A workable weaker condition is conditional
selection on pre-treatment information:
\[
    R_j\perp\theta_j
    \mid
    X_j,S_j,
\]
together with overlap
\[
    0<\pi_j(X_j,S_j)<1
\]
for every target stratum. In a finite-population design, the clean statement
is that the known sampling probability may depend on observed pre-treatment
\((X_j,S_j)\) but not on the unobserved gain. The resulting moments or
aggregate estimator must then use the known inclusion probabilities. Simple
unweighted moment matching is valid only when the validation sample is
self-weighting or when the channel is correctly modeled conditional on the
selection variables.

Useful stratification variables include occupation cluster, task family,
industry, firm size, baseline digital intensity, pre-AI task mix, geography,
and the broad score \(S_j\). The conditioning set must be fixed before
validation outcomes are observed. If a target stratum has zero probability of
validation, its effect is not design-identified; inference for that stratum is
an extrapolation or bounds exercise.

\subsection*{2.9 Sites are clusters, so forty cells may not mean forty
independent anchors}

The proposal describes roughly ten enterprise deployments yielding roughly
forty validated cells. Cells from one enterprise share technology, management,
worker composition, and implementation shocks. They therefore need not supply
forty independent channel observations.

Let the average number of validated cells per site be
\[
    \bar c=\frac mD
\]
and let \(\rho\) be an intraclass correlation for the information relevant to
channel estimation. A standard approximation is
\[
    m_{\rm eff}
    \approx
    \frac{m}{1+(\bar c-1)\rho}.
\]
The channel-estimation term is then closer to
\[
    \frac{\Lambda}{m_{\rm eff}}
\]
than to \(\Lambda/m\). The exact analysis should use site-level inclusion
probabilities, joint cell inclusion probabilities, and cluster-robust or
design-based covariance calculations. If only ten sites are independently
sampled, the site count \(D\), not just the raw cell count \(m\), governs the
external-validity and dependence problem.

\subsection*{2.10 Recommended validation design}

A defensible design is:
\begin{enumerate}[leftmargin=2em]
    \item Define \(\theta_j\) over an explicit target distribution of sites,
    workers, and tasks.
    \item Stratify candidate sites using pre-treatment \(X_d\) and the cells
    they cover.
    \item Randomly sample sites within strata with known \(q_d>0\).
    \item Randomly or probability-sample cells within selected sites, giving
    \(\pi_j=q_dr_{j\mid d}\).
    \item Randomize AI access or model tier within each selected site-cell.
    \item Estimate the channel with sampling weights and site-clustered
    covariance.
    \item Report both \(m\) and the number \(D\) of independent sites, along
    with a design-effect or effective-sample-size calculation.
    \item If convenience sites must be used, state the conditional transport
    assumption and overlap diagnostics. Do not call the resulting anchor
    randomized with respect to the target population.
\end{enumerate}

\section*{Main conclusions}

\begin{enumerate}[leftmargin=2em]
    \item The annular injectivity problem does not require a continuous map in
    \(\theta\). Finite and countably infinite supports are generically
    injectively encodable by a scalar real-valued index when the weight space
    is rich enough. For uncountable supports, geometry rather than cardinality
    is decisive.
    \item Not every exposure map is uninformative. The problem arises when the
    exposure encodes all assignment features used by the design functions.
    It is separate from the broader minimax impossibility result.
    \item \(Y_i\) is the shorthand random variable \(Y_i(W)\) in the
    design-based framework. Renaming it does not make off-truth score means
    design-computable.
    \item Of the five July 27 feedback items, the July 30 redline
    substantively addresses only item 0.4, the growing-sieve objection.
    \item AI held-out-cell total MSE has an oracle floor:
    \[
        \MSE_{\rm held}(m)=\omega^2+\Lambda/m+o(m^{-1}).
    \]
    Therefore a budget threshold must be stated relative to that floor.
    \item The design-based aggregate has an exact finite-population MSE
    formula. Under fixed total worker-level budget, adding cells reduces
    representativeness error but not the total within-cell noise component.
    \item Within-site treatment randomization is not a substitute for
    representative site or cell sampling. Conditional selection can replace
    unconditional random sampling only with pre-treatment conditioning,
    overlap, and correct weighting or transport assumptions.
\end{enumerate}

\section*{Files reviewed}

\begin{enumerate}[leftmargin=2em]
    \item \texttt{feedback\_exposure\_0727 (1).pdf}.
    \item \texttt{Exposure\_Feedback\_7\_28 (1) (1).pdf}.
    \item \texttt{park exposure 0727.pdf}.
    \item \texttt{park exposure 0730.pdf}.
    \item \texttt{gao harshaw savje et al 2026 impossibility of testing exposure
    mapping.pdf}.
    \item \texttt{Exposure\_Proof\_Review\_0615.pdf}.
    \item \texttt{feedback\_comments\_response\_0630.pdf}.
    \item \texttt{gao\_park\_exposure\_writeup\_revised.tex}.
    \item \texttt{openai\_exchange\_proposal\_v4\_for\_ext.pdf}.
    \item \texttt{ai\_project\_three\_papers\_whiteboard\_note.tex}.
\end{enumerate}

\end{document}
